\documentclass{umphdthesis}

%----------------
% Thesis metadata
%----------------
\thesistitle{Title of Thesis}
\thesisauthor{Given Names}{Surname}
\thesissupervisor{Name of Supervisor}
\thesiscosupervisor{Name of Co-supervisor(s) (if applicable)} % if applicable
\thesisdepartment{Name of Department (if applicable)} % optional
\thesisdegreetitle{Degree of Title}
\thesisdegreeyear{Year}
\thesisacademicunit{Academic Unit}
\thesisauthoryear{Year}

%----------------------
% Your packages go here
%----------------------
\usepackage{lipsum}

\begin{document}

%------------------------------
% Front matter (Roman numerals)
%------------------------------
\umphdfrontmatter

\begin{umphdacknowledgements}
\lipsum[1-2]
\end{umphdacknowledgements}

% Abstract (max 350 words!)
\begin{umphdabstract}
\lipsum[1-2]
\end{umphdabstract}

% Declaration
\begin{umphddeclaration}
\end{umphddeclaration}

% Table of Contents
\makeatletter
\@openrightfalse
\makeatother
\tableofcontents

% List of Tables and Figures (combined)
\umphdlistoffiguresandtables

% Glossary (optional)
\begin{umphdglossary}
\end{umphdglossary}

% List of Abbreviations (optional)
\begin{umphdabbreviations}
\end{umphdabbreviations}

%------------------------------
% Main matter (Arabic numerals)
%------------------------------
\umphdmainmatter

\chapter{Introduction}

The \gls{latex} system is used for \glspl{term}.
We also use \gls{acr:um}.

\section{Background}

\lipsum[1-2]

\section{Research Objectives}

\lipsum[1-2]

\section{Contributions and Originality}

\lipsum[1-2]

\chapter{Literature Review}

The practice of writing programs as documents intended for human
readers, known as literate programming, was introduced by
\cite{knuth1984} and remains influential in the design of modern
documentation tools. Its ideas were popularised in part through the
\LaTeX\ typesetting system \cite{lamport1994}, which itself emphasises
the separation of content from presentation.

\section{Deep Learning for Perception}
\subsection{Object Detection}
\lipsum[1-2]

\chapter{Methodology}
\section{Proposed Architecture}
\lipsum[1-2]

\chapter{Experiments and Results}
\section{Datasets}
\lipsum[1-2]

\chapter{Conclusion and Future Work}
\lipsum[1-2]

%------------------------------
% References (APA or MLA style)
%------------------------------
\bibliographystyle{apalike}
\bibliography{bibliography}
\addcontentsline{toc}{chapter}{References}
\clearpage

%-----------
% Appendices
%-----------
\appendix
\chapter{Additional Experimental Results}

\lipsum[1-2]

\end{document}